% Standalone section. Uses the established p=37 module and cup inputs.
\section{The cyclotomic core of a torsion lift at \(37\)}
\label{sec:p37-core}

At \(p=37\), write
\[
 H=\operatorname{Gal}(\Omega_{37}/F_\infty),\qquad
 P=\operatorname{Gal}(\Omega_{37}/F_1),\qquad
 \Gamma=P/H=\langle\gamma\rangle.
\]
The arithmetic inputs established above give
\[
 H^{\rm ab}=U\oplus T,\qquad U\simeq\Lambda^{18},
 \qquad T=T_{32}\simeq\mathbf Z_{37}.
\]
Write \(\gamma|_T=\alpha\in1+37\mathbf Z_{37}\).
The nonzero cup of \cite[Theorem~7.5, p.~20; Corollary~8.3, p.~23]{MS}, through the proved
companion comparison, implies that the central class-two
extension above \(T_{32}\) has no nonzero \(\Lambda\)-torsion.

\begin{theorem}\label{thm:p37-cyclotomic-core}
Let \(t\in H\) have its entire abelian image equal to a generator
of \(T\), and let \(N_t\) be its closed normal closure in \(H\).
Then
\[
 \operatorname{Core}_{P}(N_t)
   =\bigcap_{\eta\in\Gamma}\eta N_t\eta^{-1}=1.
\]
\end{theorem}

\begin{proof}
Choose lifts \(x_{i,0}\in H\) of an integral \(\Lambda\)-basis
\(u_1,\ldots,u_{18}\) of \(U\), and define
\[
 x_{i,n+1}=\gamma x_{i,n}\gamma^{-1}x_{i,n}^{-1}.
\]
Their abelian images are \((\gamma-1)^nu_i\).
The sequences \(x_{i,n}\) tend to \(1\): their images are
iterated commutators of increasing length in each finite
quotient of the pro-\(37\) group \(P\), and every finite
\(37\)-group is nilpotent.
Together with \(t\), they lift a convergent product
\(\mathbf Z_{37}\)-basis of \(H^{\rm ab}\).

These elements are a free pro-\(37\) basis. Indeed, the
homomorphism from the free pro-\(37\) group on the corresponding
convergent pointed set to \(H\) is an isomorphism on
abelianizations. Its compact image is Frattini-surjective,
hence all of \(H\). Freeness of \(H\) gives a continuous
section; its closed image in the source is again
Frattini-surjective and hence is the entire source.
Thus the homomorphism is an isomorphism.

Let \(P_0\) be the closed free factor on the \(x_{i,n}\).
Then
\[
 H=P_0*F(t),\qquad P_0^{\rm ab}=U.
\]
It is \(\Gamma\)-stable, since
\(\gamma x_{i,n}\gamma^{-1}=x_{i,n+1}x_{i,n}\).
This gives invariance under positive powers; closedness and
\(\gamma^{37^s-1}\to\gamma^{-1}\) give invariance under the
inverse, and then under all of \(\Gamma\).

Let \(q_0:H\to P_0\) fix \(P_0\) and kill \(t\), and for
\(n\ge0\) put
\[
 r_n=\gamma^{-n}|_{P_0}\circ q_0\circ\gamma^n.
\]
It fixes \(P_0\), sends \(t\) to an element \(w_n\in[P_0,P_0]\),
and has kernel \(\gamma^{-n}N_t\gamma^n\).
We show that the intersection of these kernels is trivial.

In the class-two Lie quotient with all brackets involving
\(T\) removed, let \(Y\) be the logarithm of \(t\).
There is a companion identity
\[
 \gamma Y=\alpha Y+C,\qquad
 C\in\widehat{\bigwedge}^{\,2}_{\mathbf Z_{37}}U.
\]
Here \(C\ne0\): otherwise \(Y\) would be a nonzero
\(\Lambda\)-torsion element above the generator of \(T\).
This also follows directly from the nonzero companion
residue, which is unchanged by altering the lift.

Use the integral \(\Lambda\)-basis of \(U\) for completed
tensor coordinates. The action of \(\gamma\) on an ordered
pair is multiplication by \(z=(1+X)(1+Y)\).
Iteration of the companion identity, followed by \(q_0\)
and then \(\gamma^{-n}\), gives the exact quadratic
coefficient of \(w_n\):
\[
 B_{n;a,b}(X,Y)
 =C_{a,b}(X,Y)\sum_{i=0}^{n-1}\alpha^{n-1-i}z^{i-n}
 =\frac{C_{a,b}(X,Y)}{z-\alpha}
       \left(1-\left(\frac{\alpha}{z}\right)^n\right).
                                                        \tag{*}
\]
The sum is a bounded integral series for every \(n\ge0\);
the fraction expression is used in a fraction field.

Suppose \(g\ne1\) is killed by every \(r_n\).
Give the Magnus variables for \(P_0\) assigned degree one
and the variable for \(t\) assigned degree two.
Let \(d\) be the first nonzero assigned degree of \(g-1\).
This degree exists by the finite-tail Magnus argument.
At degree \(d\) there are finitely many token patterns:
singletons of one of eighteen odd colors, and adjacent
pairs occupied by \(t\). The completed coefficients are
bounded series
\[
 A_{\mathcal P}(T_i:i\text{ singleton})
 \in\mathbf Z_{37}[[T_i:i\text{ singleton}]].
\]
Indeed each odd tensor factor is \(\Lambda\), and its
completed \(s\)-fold tensor product over \(\mathbf Z_{37}\)
is \(\mathbf Z_{37}[[T_1,\ldots,T_s]]\).

Every \(r_n\) fixes the odd generators and sends \(t\) to
ordinary degree at least two. Thus higher source degrees,
and higher terms of the \(w_n\), cannot contribute at
target degree \(d\). The degree-\(d\) substitution uses only
\((*)\). It is legitimate on the completed tensors:
finite-tail truncations converge uniformly on compact
\(H\), and the fixed-degree substitutions are continuous.

Choose a nonzero token pattern with the largest number
\(k\) of pairs, and call its pair set \(J\).
For each selected pair choose an ordered color pair on
which \(C\) has a nonzero coefficient. Retain the source
colors at its singleton positions. These choices are
independent because the pairs are disjoint.

For the \(d\) independent target positions put \(Z_i=1+T_i\)
and work in
\[
 K_d=\operatorname{Frac}
       \bigl(\mathbf Z_{37}[[T_1,\ldots,T_d]][1/37]\bigr).
\]
For a pair starting at \(j\), write
\[
 R_j=\frac{C_{a_j,a_{j+1}}(T_j,T_{j+1})}
              {Z_jZ_{j+1}-\alpha},
 \qquad \beta_j=\frac{\alpha}{Z_jZ_{j+1}}.
\]
For every integer \(n\ge0\), the vanishing target coefficient
becomes
\[
 \sum_{\mathcal P}A_{\mathcal P}
   \prod_{j\in J(\mathcal P)}R_j(1-\beta_j^n)=0.
\]
Expanding the products gives a finite sum of exponentials
\[
 \beta_I^n,\qquad
 \beta_I=\prod_{j\in I}\beta_j
    =\frac{\alpha^{|I|}}{\prod_{i\in V(I)}Z_i},
\]
where \(I\) is an adjacent disjoint-pair matching and
\(V(I)\) its covered vertex set.

The \(\beta_I\)'s are distinct. Algebraic independence of
the \(Z_i\)'s makes equality imply the same covered vertex
set. Such a set determines its adjacent matching uniquely:
the leftmost vertex must pair with its right neighbor,
and one then proceeds inductively. The identities at
\(n=0,\ldots,M-1\), where \(M\) is the number of distinct
exponentials, have an invertible Vandermonde matrix over
\(K_d\). Every exponential coefficient is therefore zero.

For \(I=J\), a source pattern can contribute only if its
pair set contains \(J\). A strictly larger set has zero
coefficient by maximality. Equality of pair sets and of
the already selected singleton colors leaves exactly
the chosen pattern, with coefficient
\[
 (-1)^k A_{\mathcal P}
   \prod_{j\in J}
   \frac{C_{a_j,a_{j+1}}(T_j,T_{j+1})}
             {Z_jZ_{j+1}-\alpha}.
\]
All its factors are nonzero in \(K_d\), a contradiction.
If \(k=0\), the retraction \(r_0=q_0\) already detects the
nonzero singleton coefficient.

Hence \(\bigcap_{n\ge0}\ker r_n=1\), proving the theorem.
Closedness and density of integral powers also identify
the integral intersection with the full procyclic one.
\end{proof}

\begin{corollary}\label{cor:one-lift-p37-faithfulness}
For every \(t\) as in Theorem~\ref{thm:p37-cyclotomic-core},
\[
 \Omega_{37}^{N_t}\subseteq\Omega_{37}^*
 \quad\Longleftrightarrow\quad
 \Omega_{37}=\Omega_{37}^*.
\]
\end{corollary}

\begin{proof}
Let \(K\le H\) be the actual outer-action kernel.
It is \(P\)-normal. Thus \(K\subseteq N_t\) implies
\[
 K\subseteq\operatorname{Core}_P(N_t)=1.
\]
The converse is immediate, and the displayed field
equivalence is the exact Galois correspondence.
\end{proof}

